% ══════════════════════════════════════════════════════════════════════
% MODELO DE RELATÓRIO DO PORTFÓLIO
% ══════════════════════════════════════════════════════════════════════
% Preâmbulo e capa copiados do analises/latex/principal.tex do projeto de
% gestão hospitalar, sem o que era específico dele (glosas de seção,
% macros da fase qualitativa, cores da fórmula do DEA).
%
% Para começar um relatório novo: copie este arquivo, preencha o bloco
% "DADOS DA CAPA" e troque a SEÇÃO DE EXEMPLO pelo conteúdo do projeto.
% O preâmbulo não se edita por relatório: ele é o kit.
%
% Compilar: pdflatex modelo_relatorio.tex (duas passadas; três se houver
% remissão a apêndice que vem depois de quem o cita). Ou latexmk -pdf.
% Motor: pdfLaTeX (MiKTeX ou TeX Live).

% ══════════════════════════════════════════════════════════════════════
% PREÂMBULO DO KIT: daqui até "FIM DO PREÂMBULO DO KIT", não edite
% ══════════════════════════════════════════════════════════════════════
\documentclass[11pt,a4paper]{article}

\usepackage[T1]{fontenc}
\usepackage[utf8]{inputenc}
\usepackage[brazilian]{babel}

% A mancha: 15 cm de linha, cerca de 3,2 alfabetos de Libertine 11pt. As
% figuras NÃO vivem nesta medida: ver \figwidth, mais abaixo.
\usepackage[top=2cm, bottom=2cm, left=3cm, right=3cm]{geometry}
\usepackage{graphicx}
\usepackage{subcaption}
\usepackage{booktabs}
\usepackage{array}
\usepackage{longtable}
\usepackage{float}
\usepackage{amsmath}
\usepackage{caption}
% A opção `table` liga o colortbl, que é o que dá a zebra (\rowcolors).
\usepackage[table]{xcolor}
\usepackage{enumitem}
% Capa: a curva em TikZ e os ícones vetoriais do Font Awesome.
\usepackage{tikz}
\usepackage{fontawesome5}

% ── AS FONTES: serifa no corpo, Lato em tabela e figura ───────────────
% `lato` SEM a opção `default`: o corpo não muda e o comando \lato fica
% disponível. A `libertine` entra DEPOIS da `lato`, porque as duas mexem
% em \sfdefault e a última vence. A matemática acompanha a Libertine.
\usepackage{lato}
\usepackage{libertine}
\usepackage{libertinust1math}

% ── AS CORES: ColorBrewer Set1 e a escala tirada do #0A0310 ───────────
% Os mesmos hexadecimais do estilo_portfolio.py e do theme_portfolio.R.
\definecolor{azul}{HTML}{377EB8}       % destaque principal
\definecolor{vermelho}{HTML}{E41A1C}   % destaque secundário
\definecolor{verde}{HTML}{4DAF4A}
\definecolor{roxo}{HTML}{984EA3}
\definecolor{laranja}{HTML}{FF7F00}
\definecolor{marrom}{HTML}{A65628}
\definecolor{rosa}{HTML}{F781BF}
\definecolor{tinta}{HTML}{0A0310}      % o quase-preto do kit
\definecolor{tinta2}{HTML}{403A45}     % GRAY_70: texto secundário
\definecolor{tinta3}{HTML}{5D5961}     % GRAY_60: rodapé, nota
% Versões de TEXTO das cores claras, escurecidas até 4,5:1 contra o branco
% por cor_de_texto() do estilo_portfolio.py. Vermelho e roxo já passam.
\definecolor{azulTexto}{HTML}{3579B1}
\definecolor{verdeTexto}{HTML}{3C823B}

% ── A ZEBRA, aplicada a TODA tabela sem editar tabela nenhuma ─────────
% A alternância é ancorada no PRIMEIRO \midrule de cada tabular (o fim do
% cabeçalho), e não no número da linha: assim a primeira linha de dado
% sai sempre branca, com cabeçalho de uma ou de duas linhas.
%
% ⚠ A ordem {zebraclara}{white} é contraintuitiva e está certa: o
% contador sobe ANTES de a cor ser escolhida. Invertê-la inverte a zebra.
% ⚠ O reset vai dentro de \noalign, senão "Misplaced \noalign".
\usepackage{etoolbox}
\definecolor{zebraclara}{HTML}{F4F4F4}
\newif\ifzebraprimeiromidrule
\newcommand{\zebratabela}{%
  \lato
  \global\zebraprimeiromidruletrue
  \rowcolors{1}{zebraclara}{white}%
}
\AtBeginEnvironment{tabular}{\zebratabela}
\AtBeginEnvironment{longtable}{\zebratabela}
\pretocmd{\midrule}{%
  \noalign{%
    \ifzebraprimeiromidrule
      \global\zebraprimeiromidrulefalse
      \global\rownum=0\relax
    \fi}%
}{}{\typeout{ATENCAO: o patch de \string\midrule\space para a zebra FALHOU}}

% ── A BOLINHA DE ESTADO, para coluna sim/não de tabela ────────────────
% É reforço, e não substituição: o texto fica ao lado. Os três tons têm o
% mesmo contraste (2,3:1), calculado no projeto de origem. Glifo, e não
% TikZ: \tikz dentro de \caption derruba a compilação.
\definecolor{bolinhaSim}{HTML}{6DBD6B}
\definecolor{bolinhaNao}{HTML}{F29192}
\definecolor{bolinhaEmpate}{HTML}{B1AE28}
\DeclareRobustCommand{\bolinha}[1]{\textcolor{#1}{%
  \raisebox{-0.06ex}{\scalebox{2}{$\bullet$}}}}
\DeclareRobustCommand{\marcasim}{\bolinha{bolinhaSim}\,sim}
\DeclareRobustCommand{\marcanao}{\bolinha{bolinhaNao}\,não}
\DeclareRobustCommand{\marcaempate}{\bolinha{bolinhaEmpate}\,empate}

% ── A CAIXA DE DESTAQUE ───────────────────────────────────────────────
% Para o achado principal de uma seção, uma por seção no máximo. Fio azul
% à esquerda, fundo da zebra, título em Lato. Quebra entre páginas.
%   \begin{destaque}[Título opcional] texto \end{destaque}
\usepackage[most]{tcolorbox}
\newtcolorbox{destaque}[1][Em uma frase]{%
  enhanced, breakable,
  colback=zebraclara, colframe=azul,
  boxrule=0pt, leftrule=2.5pt, arc=0pt, outer arc=0pt,
  left=10pt, right=10pt, top=8pt, bottom=8pt,
  before skip=12pt, after skip=12pt,
  % ⚠ A cor do título é `coltitle`: um \color antes de \tcbtitle não vale,
  % porque o tcolorbox reaplica coltitle (branco por padrão) por dentro.
  title={#1}, detach title, coltitle=azulTexto,
  before upper={{\lato\bfseries\tcbtitle}\par\smallskip},
}

\usepackage{needspace}
\usepackage[hidelinks]{hyperref}
\setcounter{tocdepth}{2}
\hypersetup{bookmarksdepth=3, bookmarksnumbered=true, bookmarksopen=false}

\captionsetup{font=small, labelfont=bf}
\captionsetup[sub]{font=footnotesize}

\graphicspath{{figuras/}{exemplos/}}

% Paginação: altura natural, sem esticar o espaço vertical; sem órfãs nem
% viúvas; folga mínima de composição em vez de reescrever prosa.
\raggedbottom
\widowpenalty=10000
\clubpenalty=10000
\emergencystretch=1em

% ── \figwidth: a figura é mais larga que a prosa ──────────────────────
% A figura vive na medida de 1 cm de margem no papel (19 cm em A4), e não
% na mancha de 15 cm. Definida pelo papel, e não por número copiado.
\newlength{\figwidth}
\setlength{\figwidth}{\dimexpr\paperwidth-2cm\relax}

% \imagemplena[fração]{arquivo}: a imagem centrada sobre a mancha. A caixa
% de largura declarada impede o transbordo só para a direita e não conta
% como Overfull \hbox.
\newcommand{\imagemplena}[2][1]{%
  \makebox[\linewidth][c]{\includegraphics[width=#1\figwidth]{#2}}}

% tabelaplena: a tabela que não cabe na mancha ganha a medida da figura.
% Só para tabular (caixa não quebra entre páginas).
\newsavebox{\tabelacaixa}
\newenvironment{tabelaplena}
  {\begin{lrbox}{\tabelacaixa}}
  {\end{lrbox}\makebox[\linewidth][c]{\usebox{\tabelacaixa}}}

% \subfig[largura]{arquivo}{legenda}: um painel por linha, na medida da figura.
\newcommand{\subfig}[3][0.62\figwidth]{%
  \makebox[\linewidth][c]{%
  \begin{subfigure}[t]{#1}
    \centering
    \includegraphics[width=\linewidth]{#2}
    \caption{#3}
  \end{subfigure}}\par\medskip}
% \linhasde{arquivo}: as linhas de uma tabela, escritas por script.
% ⚠ Não use \input dentro de tabular: ele não é expansível, e o
% \bottomrule seguinte derruba a compilação com "Misplaced \noalign". O
% \@@input primitivo é expansível, e o espaço depois de #1 encerra o nome.
\makeatletter
\newcommand{\linhasde}[1]{\@@input #1 }
\makeatother
% ══════════════════════════════════════════════════════════════════════
% FIM DO PREÂMBULO DO KIT
% ══════════════════════════════════════════════════════════════════════

% ── DADOS DA CAPA: preencha por relatório ─────────────────────────────
\newcommand{\RotuloCapa}{RELATÓRIO ESTATÍSTICO \textperiodcentered\ PROJETO DE PORTFÓLIO}
\newcommand{\TituloRelatorio}{Kit de Estilo do Portfólio}
\newcommand{\SubtituloRelatorio}{Modelo de relatório, com uma seção de exemplo sobre dados fictícios}
\newcommand{\DelimitacaoRelatorio}{Dados, recorte e período do estudo \\ Período: \NumAnoInicial{} a \NumAnoFinal}
\newcommand{\AutorNome}{Seu Nome}
\newcommand{\AutorPapel}{Estatístico responsável}
\newcommand{\AutorLinkedIn}{https://www.linkedin.com/in/seu-perfil/}
\newcommand{\AutorGitHub}{https://github.com/seu-usuario}
\newcommand{\LocalCapa}{São Paulo -- SP}
\newcommand{\DataCapa}{Novembro de 2026}

% ── OS NÚMEROS DO RELATÓRIO: gerados por script, nunca digitados ──────
% Cada \Num... vem de um arquivo escrito pelo script que calcula o número.
% Ver GUIA_ESTILO.md, "A regra dos números".
\input{exemplos/numeros_exemplo.tex}

\title{\TituloRelatorio}

\begin{document}

% ══════════════════════════════ CAPA ══════════════════════════════════
% Azul como destaque principal e vermelho como secundário. As quatro cores
% principais da paleta aparecem juntas em um só lugar: os traços ao pé da
% curva. Nenhum outro elemento da capa usa verde ou roxo.
\begin{titlepage}
  \centering
  \vspace*{0.5cm}

  {\color{tinta}\rule{\linewidth}{1.5pt}} \\[0.4cm]
  {\small\bfseries\color{tinta2} \RotuloCapa} \\[0.2cm]
  {\color{tinta}\rule{\linewidth}{0.5pt}}

  \vspace{1.6cm}

  % A curva ecoa a série temporal; ela e os quatro traços cobrem o mesmo vão.
  \begin{tikzpicture}
    \draw[line width=1.6pt, azul, smooth]
      plot coordinates {(0.2,0.35) (1.7,0.62) (3.2,0.48) (4.7,0.95)
                        (6.2,1.32) (7.7,1.12) (9.2,1.58) (10.7,1.42)};
    \foreach \x/\c in {0.2/verde, 2.9/vermelho, 5.6/azul, 8.3/roxo}
      \fill[\c] (\x,-0.28) rectangle ++(2.4,0.07);
  \end{tikzpicture}

  \vspace{1.2cm}

  \begin{minipage}{0.94\textwidth}
    \centering
    {\Huge\bfseries\color{azul} \TituloRelatorio \par}
    \vspace{0.5cm}
    {\LARGE\bfseries\color{vermelho} \SubtituloRelatorio \par}
  \end{minipage}

  \vspace{1.1cm}

  % Dois a quatro ícones, um por eixo do estudo; troque pelos do projeto.
  % ⚠ Use os nomes sem sufixo `Alt` (\faMapMarker, \faCalendarDay): o
  % fontawesome5 do MiKTeX não define os `...Alt`.
  {\Large
   \textcolor{azul}{\faDatabase}\hspace{1.1cm}
   \textcolor{vermelho}{\faChartLine}\hspace{1.1cm}
   \textcolor{azul}{\faCalculator}\hspace{1.1cm}
   \textcolor{vermelho}{\faCalendarDay}\par}

  \vspace{1.0cm}

  {\large\itshape\color{tinta2} \DelimitacaoRelatorio \par}

  \vspace{2.2cm}

  {\large\bfseries\color{azul} Autoria} \\[0.6cm]

  \begin{minipage}{0.85\textwidth}
    \centering
    {\large\bfseries \AutorNome} \\
    {\normalsize\color{tinta2} \AutorPapel} \\[0.15cm]
    {\footnotesize
     \href{\AutorLinkedIn}{\color{vermelho}\textbf{LinkedIn}} \ \textperiodcentered\ \
     \href{\AutorGitHub}{\color{vermelho}\textbf{GitHub}}}
  \end{minipage}

  \vfill

  {\color{tinta2}\rule{4cm}{0.5pt}} \\[0.3cm]
  {\small\color{tinta2} \LocalCapa \\ \bfseries \DataCapa}
\end{titlepage}

% Em relatório longo, ligue o sumário:
% \tableofcontents
% \clearpage

% ══════════════════════════════════════════════════════════════════════
% SEÇÃO DE EXEMPLO: troque pelo conteúdo do projeto
% ══════════════════════════════════════════════════════════════════════
% Todo número abaixo é uma macro \Num... escrita por
% exemplos/exemplo_figura.py. Nenhum foi digitado.
\section{A espera caiu onde a triagem mudou}

Esta seção mostra o kit em uso sobre \textbf{dados fictícios}: seis
unidades de atendimento, com o tempo médio de espera para a primeira
consulta medido de \NumAnoInicial{} a \NumAnoFinal. Em \NumAnoMudanca, a \NumFoco{} adotou uma
nova triagem; as outras cinco não mudaram nada.

\begin{destaque}[O achado, em uma frase]
  A \NumFoco{} passou de \NumEsperaAntes{} para \NumEsperaDepois{} dias de
  espera entre \NumAnoAntes{} e \NumAnoFinal, uma queda de
  \NumQuedaFoco\%. No mesmo período, nenhuma das outras unidades variou
  mais de \NumVarOutrasMax\%.
\end{destaque}

A Figura~\ref{fig:exemplo} segue as regras do kit: o título afirma o
achado, a série que importa tem a única cor da figura e as demais ficam no
neutro, o nome vai na ponta da linha em vez de numa legenda e a fonte dos
dados fica no rodapé. O eixo começa em zero, e a figura ocupa a medida de
imagem (\verb|\imagemplena|), mais larga que a linha do texto.

\begin{figure}[H]
  \centering
  \imagemplena[0.86]{exemplo_figura}
  \caption{Tempo médio de espera para a primeira consulta, em dias, nas
    seis unidades, de \NumAnoInicial{} a \NumAnoFinal. A \NumFoco{} (em vermelho) caiu
    \NumQuedaFocoExata\% depois da nova triagem; a linha tracejada marca a
    mudança. Dados fictícios, gerados com semente fixa pelo script de
    exemplo do kit.}
  \label{fig:exemplo}
\end{figure}

A Tabela~\ref{tab:exemplo} traz os mesmos números para quem precisa do
valor exato. Ela está ordenada pela espera em \NumAnoFinal, e não pelo nome: a
ordem já responde quem está mais perto da meta de \NumMetaDias{} dias.
Das \NumUnidades{} unidades, \NumUnidadesNaMeta{} estão dentro dela.

\begin{table}[H]
  \centering
  \caption{Espera média antes e depois da mudança, por unidade (dias).
    A coluna da meta usa a bolinha como reforço do texto, e não no lugar
    dele. Dados fictícios.}
  \label{tab:exemplo}
  \small
  \begin{tabular}{@{}lrrrl@{}}
    \toprule
    Unidade & \NumAnoAntes & \NumAnoFinal & Variação & Meta de \NumMetaDias{} dias \\
    \midrule
    \linhasde{exemplos/tabela_exemplo.tex}
    \bottomrule
  \end{tabular}
\end{table}

\paragraph{Limitações.} Os dados são fictícios e a comparação é
descritiva: uma queda que coincide com a mudança não prova que a mudança a
causou. Num projeto real, esta é a subseção em que se diz o que a análise
não autoriza afirmar.

\end{document}
